% paper.tex —— "Geometry by Construction: A Prior-Geometry-Constrained Information Bottleneck"
% 基于 paper/OPB.txt（分类）、paper/OPB-R.txt（回归）、paper/theory.tex（理论）撰写；
% 结构参照 paper/template.tex（ICLR 2027 投稿格式）。本仓库未随附 iclr2027_conference.sty
% 与 math_commands.tex，故 preamble 自带最小等价设置（拿到样式文件后替换即可）。
% 当前包含：题目、Intro、方法（2.1 背景与记号 / 2.2 关键理论 / 2.3 模型变体）、
% 附录（其余理论结果）。实验与相关工作为占位，后续填充。

\documentclass{article}

% ---------- 最小等效 ICLR 样式（原模板为 \usepackage{iclr2027_conference,times}） ----------
\usepackage[T1]{fontenc}
\usepackage{times}
\usepackage{amsmath,amssymb,amsthm}
\usepackage{natbib}
\usepackage{hyperref}
\usepackage{url}

\theoremstyle{plain}
\newtheorem{proposition}{Proposition}[section]
\newtheorem{lemma}{Lemma}[section]
\newtheorem{corollary}{Corollary}[section]
\theoremstyle{remark}
\newtheorem{remark}{Remark}[section]

\newcommand{\R}{\mathbb{R}}
\newcommand{\E}{\mathbb{E}}
\DeclareMathOperator{\KL}{KL}
\DeclareMathOperator{\CE}{CE}

\title{Geometry by Construction: A Prior-Geometry-Constrained Information Bottleneck}

% 匿名投稿：作者信息隐藏（同 template.tex 约定）
\author{Anonymous authors}
%\iclrfinalcopy % 相机版时取消注释

\begin{document}

\maketitle

% =====================================================================
\section{Introduction}
% =====================================================================

The information bottleneck (IB) principle \citep{tishby2000information} seeks a
representation $Z$ of an input $X$ that retains the information necessary to predict a
label $Y$ while discarding the rest. Because the mutual information terms are
intractable for deep networks, the variational information bottleneck (VIB)
\citep{alemi2017deep} bounds the compression term $I(X;Z)$ with a fixed marginal prior
$\mathcal{N}(0, I)$; the conditional entropy bottleneck (CEB) \citep{fischer2020conditional}
replaces the marginal prior with a trainable label-conditional backward encoder
$b(z|y)$ and thereby sharpens the bound from the full rate to the \emph{residual}
information $I(X;Z|Y)$, certified by the gap
$\mathrm{ReX} = \E[\KL(q(z|x)\|b(z|y))] \ge I(X;Z|Y)$.
Nonlinear IB \citep{kolchinsky2017nonlinear} and the multivariate variational
formulation of \citep{abdelaleem2025deep} complete the deep-variational lineage.
The certificate is the central object of this paper, and its failure modes are our
starting point.

\textbf{Two failure channels.} In the deterministic scenario $Y = f(X)$ --- the regime
of classification benchmarks --- the conditional bottleneck degenerates in two
independent ways. First, \emph{the knob is dead}: as shown by \citet{kolchinsky2019caveats}
(their Caveat~1), every $\gamma > 0$ of the CEB objective selects the same corner of the
information plane, $I(Y;Z) = H(Y)$ with $I(X;Z|Y) = 0$, because the CEB weight is
$\beta_{\mathrm{IB}} = 1 + \gamma > 1$ and the deterministic IB curve is linear. We
restate this in our notation in Proposition~\ref{prop:dead} and cite the proof.
Second, and less recognized, \emph{the certificate is blind}: on the family
$T_\alpha = B_\alpha \cdot Y$ that interpolates between the collapsed code and the exact
copy of the label, $\mathrm{ReX}$ reads zero at \emph{every} $\alpha$, and the objective
trades label information against certified residual at a fixed exchange rate of
$1{:}\beta$ nats. We prove this channel in Proposition~\ref{prop:flat}; its consequence
is that the certificate can be tight while the representation retains all, part, or none
of the label --- tightness of the bound says nothing about compression of the deployed
code.

\textbf{The missing ingredient.} Both channels share a single root cause: CEB's
backward encoder is an \emph{unconstrained trainable map}. It can eliminate the gap on
its own by chasing the posterior ($b_\theta \to q(z|y)$ with the encoder held fixed),
which is what blinds the certificate; and because nothing constrains the geometry of
its class means, the class-conditional priors $\{b(z|y=k)\}$ carry no explicit
relationship to one another --- two class means may sit arbitrarily close, in arbitrary
directions, and the objective never distinguishes these cases. CEB therefore compresses
\emph{against a reference with no declared structure}: the target of the matching is
whatever the optimizer finds convenient, and the certificate measures proximity to that
moving target rather than a property of the representation itself. What a
representation should look like when it is properly compressed --- class means
separated, ordered, or laid out on a frame --- is exactly what the objective leaves
unspecified. The empirical literature points the same way: deep classifiers trained
to the terminal phase collapse their class means onto an orthogonal simplex frame
\citep{papyan2020prevalence} --- a separation geometry that \emph{emerges} late in
training as a byproduct, rather than being imposed by the objective. Our framework
turns exactly this structure into a declared prior. This is the gap the present paper
closes.

\textbf{A prior-geometry-constrained information bottleneck.} We propose a framework in
which the label-conditional prior is \emph{confined to a geometric family by
construction}. The general principle is simple: the prior $p(z|y)$ must encode, through
its geometry alone, the structure of the label space, and it must do so in a way the
optimizer cannot undo. For discrete labels the structure to encode is \emph{class
equality} --- classes are either identical or distinct, nothing else --- and for
continuous labels it is \emph{order and numerical distance}. Concretely, the prior is
computed at every forward pass by a stateless geometric operation on the prior
encoder's output: the resulting prior family cannot merge two labels, cannot collapse
to a point, and cannot inflate its covariance, because these degenerations are
excluded by geometry rather than by the optimizer's good behavior. The sample-level
posterior, the single-KL objective, and the residual-certificate structure of CEB are
kept unchanged, so the entire analysis of the certificate transfers pointwise.
We call this family of objectives a \emph{prior-geometry-constrained information
bottleneck} (PG-IB). Two instantiations realize the principle.

For classification, the prior encoder is evaluated on \emph{all} class one-hot vectors
at each step, the resulting class-mean matrix is orthonormalized by a reduced QR
factorization with a fixed sign convention, and the prior means become the scaled
orthogonal frame $a\,q_{\theta,k}$: every pair of distinct classes sits at distance
$\sqrt{2}\,a$, with a class-mean Gram matrix equal to $a^2 I_K$ \emph{by construction}.
We call this variant the \emph{orthogonal-prior bottleneck} (OPB). For regression, the
prior is a line: the prior means are $\rho\,\tilde{y}\,u$ for a unit direction $u$
normalized from a trainable vector at every step, so that latent distances are exactly
$\rho$ times standardized label distances. We call this variant the \emph{isometric
regression prior} (OPB-R). Both variants are stateless --- no EMA, no prototypes, no
per-batch statistics --- and both keep a single KL term.

The geometric prior earns its keep analytically. We prove three guarantees for the
classification variant: (i) the stateless construction \emph{preserves the CEB
certificate} --- the residual bound $I(X;Z|Y)$ remains valid at every parameter value,
with a caveat that a batch-dependent QR would break it
(Proposition~\ref{prop:opb-certificate}); (ii) a class-independent posterior mean pays a
\emph{positive compression cost} of at least $a^2(1-\sum_k \pi_k^2)/(2\tau^2)$ --- the
VIB corner $\mu \equiv 0$ is a penalized configuration rather than a free optimum
(Proposition~\ref{prop:opb-collapse}); and (iii) whenever the class-wise anchor mismatch
is small, the posterior class centers inherit the separation of the frame:
$\|m_i - m_j\| \ge \sqrt{2}\,a - \sqrt{2\tau^2 D_i} - \sqrt{2\tau^2 D_j}$
(Proposition~\ref{prop:opb-separation}). The last is the transfer guarantee that turns
the prior's geometry into a property of the deployed representation. Further results ---
the KL decomposition, the anchor-energy classifier equivalence, a Fano-type
label-information lower bound, and the strict-convexity of the $\beta$-sweep --- are
collected in the appendix.

We evaluate the framework empirically across nine datasets and four backbones, and we
verify the mechanism itself rather than only end-to-end accuracy: we measure the
geometry of the learned priors and of the posterior class centers, test whether a
parameter-free nearest-anchor rule can replace the trained classifier, and ablate the
classifier with an anchor-energy rule that cannot bypass the frame. The experiments
show that the orthogonal geometry transfers from the prior to the posterior, that the
anchor scale directly controls the separation and cross-seed consistency of the
representation, and that prediction quality survives when the classifier is forced to
use the frame alone.

Our contributions are: (1) the prior-geometry-constrained IB framework, in which the
label-conditional prior encodes the structure of the label space by construction; (2)
the OPB and OPB-R variants --- a stateless orthogonal frame for classification and a
strictly isometric axis for regression --- with a single KL and a preserved CEB
certificate; (3) a theory of the two failure channels (one cited, one proven here) and
the three guarantees above; (4) a mechanism-level empirical validation that the
constructed geometry is transferred to the posterior and is sufficient for prediction.

% =====================================================================
\section{Method}
% =====================================================================

\subsection{Background and setup}
\label{sec:background}

Let $X$ be the input and $Y$ the label with $K$ classes; the encoder induces the Markov
chain $Y - X - Z$. Following \citet{fischer2020conditional}, the variational CEB
objective is
\begin{equation}
\mathrm{VCEB}
= \min_{e,b,c}\ \bigl\langle \log e(z|x)\bigr\rangle
- \bigl\langle \log b(z|y)\bigr\rangle
- \gamma\,\bigl\langle \log c(y|z)\bigr\rangle ,
\label{eq:vceb}
\end{equation}
where $e(z|x)$ is the stochastic encoder, $b(z|y)$ the backward encoder, $c(y|z)$ the
classifier, and the expectations are over $p(x,y)\,e(z|x)$. We write the encoder as the
diagonal Gaussian posterior $q_\phi(z|x) = \mathcal{N}(\mu_\phi(x),
\operatorname{diag}(\sigma_\phi^2(x)))$ and use the loss convention
$\CE + \beta\,\KL$, with $\beta = 1/\gamma$; at evaluation the deterministic code
$z = \mu_\phi(x)$ is used. The central object is the residual certificate: for any
conditional prior $r(z|y)$,
\begin{equation}
\E\bigl[\KL(q_\phi(z|x)\,\|\,r(z|y))\bigr]
= I(X;Z|Y) + \E_y\bigl[\KL(q_\phi(z|y)\,\|\,r(z|y))\bigr]
\;\ge\; I(X;Z|Y) ,
\label{eq:gap}
\end{equation}
by Fischer's chain-rule identity (their Eqs.~4 and 9--11, 20). CEB instantiates
$r(z|y) = b_\theta(z|y)$ with a trainable linear map of the one-hot label; VIB fixes
$r(z|y) = \mathcal{N}(0, I)$ for all $y$. Our framework replaces $r$ with a geometric
prior family, defined in Section~\ref{sec:construction}.

\subsection{Key theory: two failure channels and the guarantees of a geometric prior}
\label{sec:theory}

\subsubsection{Channel 1 (dead knob): every $\gamma$ selects the same corner}

\begin{proposition}[Dead knob, following Caveat~1 of \citet{kolchinsky2019caveats}]
\label{prop:dead}
Assume $Y = f(X)$. For every $\gamma > 0$, the CEB objective
$L_\gamma(Z) := I(X;Z|Y) - \gamma I(Y;Z)$ is minimized exactly on the corner family
$\{Z : I(Y;Z) = H(Y),\; I(X;Z|Y) = 0\}$, and the map
$\gamma \mapsto \bigl(I(X;Z^*|Y),\, I(Y;Z^*)\bigr)$ is constant. Consequently no
interior point of the IB curve is ever a Lagrangian minimizer of CEB, over the entire
parameter range.
\end{proposition}

\begin{proof}[Proof sketch (see \citet{kolchinsky2019caveats} for the full argument)]
By the chain rule, $L_\gamma(Z) = I(X;Z) - (1+\gamma)I(Y;Z)$; the data-processing
inequality gives $L_\gamma \ge -\gamma H(Y)$, attained exactly at the corner
$Z = Y$. In the Caveats paper's maximizing convention this is their Eq.~7 with
$\beta' = 1/(1+\gamma) \in (0,1)$, the range in which every maximizer is pinned at the
copy corner.
\end{proof}

The corner is a single point of the information plane and an entire equivalence class
of representations; the objective is exactly indifferent within the class, and neither
the objective nor the certificate records which member the optimizer reaches. In the
deterministic scenario this corner \emph{is} the MNI point, so CEB's declared target
and the degenerate attractor coincide; the $\gamma$ knob acts on the optimization
trajectory, not on the optimum.

\subsubsection{Channel 2 (blind certificate): tight at every retention level, with a $1{:}\beta$ exchange rate}

\begin{proposition}[Blind certificate and the $1{:}\beta$ exchange rate]
\label{prop:flat}
On the manifold $T_\alpha = B_\alpha \cdot Y$ of \citet{kolchinsky2019caveats}
(their Eq.~6), $I(X;T_\alpha|Y) = 0$ for every $\alpha$, and with the optimal backward
encoder and classifier,
\begin{equation}
\mathrm{ReX}(T_\alpha) = 0,
\qquad
\CE(T_\alpha) = H(Y) - I(Y;T_\alpha) =: \delta_\alpha,
\qquad
\mathrm{VCEB}(T_\alpha) = \gamma\,\delta_\alpha .
\label{eq:vceb-flat}
\end{equation}
Consequently: \emph{(a)} the certificate is blind to label retention --- $\mathrm{ReX}$
reads $0$ from the total collapse ($\alpha = 0$) to the full copy ($\alpha = 1$), so
tightness is compatible with retaining all, part, or none of $Y$; \emph{(b)} the
objective exchanges $1$~nat of label information against $\beta$~nats of certified
residual --- on $T_\alpha$, $L = \CE + \beta\,\mathrm{ReX} = \delta_\alpha$, so in the
heavy-regularization regime the entire label axis, up to $H(Y)$ nats, is purchasable
for $H(Y)/\beta$ nats of certified compression.
\end{proposition}

\begin{proof}
$T_\alpha$ is a function of $Y$ alone, hence $T_\alpha \perp X \mid Y$ and
$I(X;T_\alpha|Y) = 0$. With the true backward encoder ($b = p(z|y)$) the certificate
is tight: $\mathrm{ReX} = I(X;T_\alpha|Y) = 0$. The optimal classifier attains
$\langle \log c(y|z)\rangle = -H(Y|T_\alpha)$, i.e.\ $\CE = H(Y|T_\alpha) = \delta_\alpha$,
and $\mathrm{VCEB} = \gamma\,\delta_\alpha$ follows from Eq.~\eqref{eq:vceb}.
Part~(b) is the definition of the objective evaluated at $\CE = \delta_\alpha$.
\end{proof}

\begin{remark}[The flatness is an exchange rate, not literal flatness]
\label{rem:flatness}
The prediction term is weaker than the compression term by exactly the factor $\beta$,
so \emph{in units of the compression term} the whole manifold lies within
$H(Y)/\beta$ of optimal. Under a scale-invariant optimizer (Adam) only the ratio
matters. The operative statement is that the label axis is \emph{cheap and unmetered}:
dropping $\delta$ nats of label information buys $\delta/\beta$ nats of certified
residual, and no value of the certificate can expose the trade.
\end{remark}

The two channels locate the disease precisely: the reference of the matching --- CEB's
backward encoder --- is an unconstrained trainable map, so it can eliminate the gap
alone (Channel~2) and its class geometry carries no declared structure (the root cause
of both). The construction below removes the mechanism by replacing the reference with
a geometric prior that cannot chase the posterior and cannot merge classes.

\subsubsection{A geometric prior by construction}
\label{sec:construction}

For classification, let $K \le d$. The prior encoder is evaluated on all $K$ one-hot
labels at the current parameter value $\theta$, producing the raw class-mean matrix
\begin{equation}
M_\theta = [m_\theta(1), \ldots, m_\theta(K)] \in \R^{d \times K} .
\label{eq:opb-raw}
\end{equation}
Assume $M_\theta$ has full column rank and take its reduced QR factorization with a
fixed positive-diagonal sign convention,
\begin{equation}
M_\theta = Q_\theta R_\theta, \qquad Q_\theta^{\top} Q_\theta = I_K .
\label{eq:opb-qr}
\end{equation}
The geometric prior replaces the raw means by the scaled orthogonal frame
\begin{equation}
p_\theta^{\mathrm{OPB}}(z|y=k)
= \mathcal{N}\bigl(a\,q_{\theta,k},\, \tau^2 I_d\bigr),
\qquad q_{\theta,k} := Q_\theta[:,k],
\label{eq:opb-prior}
\end{equation}
where $a \ge 0$ is the anchor scale and $\tau^2 > 0$ the shared prior variance. The QR
factor is recomputed from the current prior encoder at every training step: there is
no EMA prototype, no persistent anchor variable, and, for fixed $\theta$, $Q_\theta$
depends only on the class set --- not on the input $x$, and not on the batch.

\begin{proposition}[Orthogonal class-prior geometry by construction]
\label{prop:opb-frame}
Under Eqs.~\eqref{eq:opb-qr}--\eqref{eq:opb-prior}, the class means
$a_k = a\,q_{\theta,k}$ satisfy
\begin{equation}
a_i^{\top} a_j = a^2 \delta_{ij},
\qquad
\|a_k\| = a,
\qquad
\|a_i - a_j\|^2 = 2a^2 \ \ (i \ne j) .
\label{eq:opb-geometry}
\end{equation}
Equivalently, the class-mean Gram matrix is $a^2 I_K$: the prior supplies a
class-separated spherical code at every training step, whatever the optimizer does.
\end{proposition}

\begin{proof}
Eq.~\eqref{eq:opb-geometry} follows from $Q_\theta^{\top}Q_\theta = I_K$ after scaling
by $a$; in particular
$\|a_i - a_j\|^2 = \|a_i\|^2 + \|a_j\|^2 - 2a_i^{\top}a_j = 2a^2$.
\end{proof}

The single objective of the variant is
\begin{equation}
\mathcal{L}_{\mathrm{OPB}}
= \E_{(x,y)}\Bigl[
\CE\bigl(c_\psi(z), y\bigr)
+ \beta\, \KL\bigl(q_\phi(z|x) \,\|\, p_\theta^{\mathrm{OPB}}(z|y)\bigr)
\Bigr],
\qquad z \sim q_\phi(z|x) .
\label{eq:opb-objective}
\end{equation}
The defining change from CEB is the prior of Eq.~\eqref{eq:opb-prior}, not a second
latent variable or a second KL term.

\subsubsection{Three guarantees}

\begin{proposition}[The stateless prior preserves the CEB certificate]
\label{prop:opb-certificate}
For every fixed parameter value $\theta$ for which Eq.~\eqref{eq:opb-qr} is defined,
$p_\theta^{\mathrm{OPB}}(z|y)$ is a conditional prior independent of $x$, and the gap
identity of Eq.~\eqref{eq:gap} applies pointwise in $\theta$:
\begin{equation}
\E\Bigl[ \KL\bigl(q_\phi(z|x) \,\|\, p_\theta^{\mathrm{OPB}}(z|y)\bigr) \Bigr]
= I_{q_\phi}(X;Z|Y)
+ \E_Y\Bigl[ \KL\bigl(q_\phi(z|Y) \,\|\, p_\theta^{\mathrm{OPB}}(z|Y)\bigr) \Bigr]
\;\ge\; I_{q_\phi}(X;Z|Y) .
\label{eq:opb-certificate}
\end{equation}
The fact that $Q_\theta$ is recomputed after parameter updates does not invalidate
this identity: at each fixed $\theta$ it remains a function of the class set only.
If $Q$ were computed from the current input batch, this conditional-prior argument
would no longer apply without a new analysis.
\end{proposition}

\begin{proof}
For fixed $\theta$, apply the gap identity of Eq.~\eqref{eq:gap} to the conditional
distribution $p_\theta^{\mathrm{OPB}}(z|y)$ and the posterior $q_\phi$.
\end{proof}

\begin{proposition}[Positive rate cost of class-independent means]
\label{prop:opb-collapse}
Let $\pi_k = \Pr(Y = k)$ and suppose the posterior mean is class-independent,
$\E[\mu_\phi(X) | Y = k] = c$ for every $k$. Then the mean-mismatch part of the rate
satisfies
\begin{equation}
\E\Bigl[ \tfrac{1}{2\tau^2} \|\mu_\phi(X) - a\,q_{\theta,Y}\|^2 \Bigr]
\;\ge\; \frac{a^2}{2\tau^2}\Bigl( 1 - \sum_{k=1}^{K} \pi_k^2 \Bigr) ,
\label{eq:opb-collapse}
\end{equation}
with equality at the optimal constant mean $c = a \sum_k \pi_k q_{\theta,k}$. For
uniform classes the bound is $a^2(1 - 1/K)/(2\tau^2)$. Unlike a zero-mean marginal
prior, the geometric prior assigns every class-independent posterior a strictly
positive mean-matching cost whenever $a > 0$ and two classes have positive
probability.
\end{proposition}

\begin{proof}
By Jensen's inequality the left side is at least
$\tfrac{1}{2\tau^2}\sum_k \pi_k \|c - a q_{\theta,k}\|^2
= \tfrac{1}{2\tau^2}\bigl(\|c\|^2 - 2a\,c^{\top}\!\sum_k\pi_k q_{\theta,k} + a^2\bigr)$.
Minimizing over $c$ gives $c = a\sum_k\pi_k q_{\theta,k}$; by orthonormality
$\|\sum_k\pi_k q_{\theta,k}\|^2 = \sum_k \pi_k^2$, yielding Eq.~\eqref{eq:opb-collapse}.
\end{proof}

\begin{proposition}[KL-controlled separation of posterior class centers]
\label{prop:opb-separation}
Let $m_k := \E[\mu_\phi(X) | Y = k]$ be the posterior class centers and
\begin{equation}
D_k := \frac{1}{2\tau^2}\,
\E\Bigl[ \|\mu_\phi(X) - a\,q_{\theta,k}\|^2 \;\Big|\; Y = k \Bigr]
\label{eq:opb-class-rate}
\end{equation}
the class-wise anchor mismatch. Then for every pair $i \ne j$,
\begin{equation}
\|m_i - m_j\| \;\ge\; \sqrt{2}\,a
- \sqrt{2\tau^2 D_i} - \sqrt{2\tau^2 D_j} .
\label{eq:opb-separation}
\end{equation}
In particular, sufficiently small class-wise anchor mismatch guarantees non-zero
separation of the posterior class centers: the frame geometry transfers to the
deployed representation.
\end{proposition}

\begin{proof}
Jensen gives $\|m_k - a q_{\theta,k}\| \le \sqrt{2\tau^2 D_k}$. Apply the triangle
inequality to
$m_i - m_j = (m_i - a q_{\theta,i}) + (a q_{\theta,i} - a q_{\theta,j})
+ (a q_{\theta,j} - m_j)$ and use Proposition~\ref{prop:opb-frame}.
\end{proof}

Propositions~\ref{prop:opb-certificate}--\ref{prop:opb-separation} are the core of the
theory. The remaining results --- the pointwise KL decomposition, the anchor-energy
classifier equivalence (which removes the free classifier scale identified in
Appendix~\ref{app:sweep}), a Fano-type label-information lower bound, the strict
convexity of the $\beta$-sweep, and the scope and limitations of the framework --- are
given in the appendix.

\subsection{Model variants}
\label{sec:variants}

The framework has two concrete instantiations, one per label-space structure. Both keep
the CEB posterior, the single-KL objective, and the residual certificate unchanged;
they differ only in the geometric family that constrains the prior.

\subsubsection{OPB: an orthogonal class frame (classification)}
\label{sec:opb}

\emph{Posterior.} The encoder computes $h = f_{\mathrm{enc}}(x)$; two heads produce the
diagonal posterior $q_\phi(z|x) = \mathcal{N}(\mu_\phi(x),
\operatorname{diag}(\sigma_\phi^2(x)))$; the code $z$ is sampled at training and set to
$z = \mu_\phi(x)$ at evaluation. The classifier consumes $z$.

\emph{Geometric prior.} The prior encoder is a single unactivated linear layer mapping
one-hot labels to $(m_\theta(k), \log \sigma_{\theta}^2(k))$. At every forward pass the
\emph{full} class one-hot matrix is passed through it, the mean matrix
$M_\theta \in \R^{d\times K}$ of Eq.~\eqref{eq:opb-raw} is orthonormalized by reduced
QR with a positive-diagonal sign convention (Eq.~\eqref{eq:opb-qr}, preventing column
sign flips across steps), and the prior becomes the scaled frame of
Eq.~\eqref{eq:opb-prior}. The loss is Eq.~\eqref{eq:opb-objective}: a single KL.

\emph{Statelessness.} $Q_\theta$ is an immediate function of the current prior encoder:
the next step recomputes it. There is no EMA, no prototype table, and no dependence on
the batch composition --- the property that makes
Proposition~\ref{prop:opb-certificate} exact.

\emph{Initialization.} The variance network and the posterior logvar head are
zero-initialized ($\sigma^2 = \tau^2 = 1$ at initialization, so training starts at
$\KL \approx \tfrac{1}{2\tau^2} a^2$); the mean block of the prior encoder is
identity-initialized, so that the initial frame $Q_\theta = [e_1, \dots, e_K]$ is full
column rank and the QR backward is stable.

\emph{Theory--code layering.} The theory states the prior with the shared variance
$\tau^2 I_d$; the released implementation uses the per-class learnable
$\sigma_{\theta}^2(k)$ output by the prior encoder (the diagonal of the variance
block). All propositions are stated for the shared-$\tau^2$ prior; the learnable
variance is a strictly more flexible family within which the shared case is the
initialization, and the experiments use the learnable version. The QR backward is
well-defined while $M_\theta$ has full column rank, which the identity initialization
and the sign convention maintain throughout the runs we report.

\subsubsection{OPB-R: an isometric regression axis (regression)}
\label{sec:opbr}

\emph{Setup.} For a one-dimensional continuous label, standardize once on the training
set: $\tilde y = (y - y_{\mathrm{center}})/y_{\mathrm{scale}}$, with the statistics
frozen for validation and test. Let $W \in \R^{d \times 1}$ be a trainable direction
and $u = W/\|W\|$ its per-step normalization. The isometric prior is
\begin{equation}
p^{\mathrm{OPB-R}}(z|y)
= \mathcal{N}\bigl(\rho\, \tilde y\, u,\; \tau^2 I_d\bigr),
\qquad
\| \mu_p(y_i) - \mu_p(y_j) \| = \rho\, |\tilde y_i - \tilde y_j| ,
\label{eq:opbr-prior}
\end{equation}
where $\rho > 0$ is the isometric scale. The isometry is exact and global: any nonzero
linear map of a scalar label is isometric up to its norm, and the normalization fixes
that norm to $\rho$ --- unlike fixed random-feature anchors, whose pairwise distances
grow nonlinearly and saturate, the isometric axis preserves label distances at every
scale (with the price that anchor norms grow with $|\tilde y|$; a line and a fixed-radius
sphere intersect in at most two points, so isometry and equal norms cannot both hold).
All prior means lie on one line; for $y_i < y_j < y_k$ the distances add,
$\|\mu_p(y_i)-\mu_p(y_k)\| = \|\mu_p(y_i)-\mu_p(y_j)\| + \|\mu_p(y_j)-\mu_p(y_k)\|$:
order and magnitude of label differences are encoded exactly.

\emph{Axis choices.} Three variants of the direction: (A) \emph{fixed axis},
$u = [1, 0, \dots, 0]$, with no trainable prior direction --- the theoretically
cleanest ablation, at the price of prespecifying a latent coordinate; (B)
\emph{trainable normalized axis}, $u = W/\|W\|$ with $W$ trained and normalized at
every step --- our main variant, which keeps the scale $\rho$ exact while allowing the
model to choose the axis; (C) \emph{unconstrained linear axis}, $\mu_p(y) = W \tilde y$
--- the simplest ablation, still isometric by linearity but with a free scale that can
shrink and weaken the geometry. We recommend (B), use (A) to test whether the trainable
rotation carries the benefit, and keep (C) as a diagnostic. The main variant uses no
bias: an affine offset cancels in pairwise distances but is a free extra degree of
matching, and $\tilde y$ is already centered.

\emph{Prediction head.} Two choices: the \emph{tied projection head}, which predicts on
the same axis the KL uses,
$\hat{\tilde y} = u^{\top} z / \rho$, then rescales
$\hat y = y_{\mathrm{center}} + y_{\mathrm{scale}}\, \hat{\tilde y}$ --- geometrically
consistent, with no free direction or scale that could bypass the prior; and the
\emph{ordinary regression head} $\hat y = \mathrm{regressor}(z)$, kept for a fair
comparison against CEB. The tied head is the recommended configuration; the ordinary
head is the comparability ablation (and is the one currently used by the released
implementation). The total objective is
$\mathcal{L}_{\mathrm{OPB-R}} = L_{\mathrm{reg}}(\hat y, y)
+ \beta\,\KL\bigl(q_\phi(z|x)\,\|\,p^{\mathrm{OPB-R}}(z|y)\bigr)$ --- a single KL, no
EMA, no label bucketing, no batch-level QR.

\emph{Relation to the classification variant.} OPB expresses the discrete structure of
the label space --- every pair of distinct classes at the same distance $\sqrt2 a$;
OPB-R expresses the continuous structure --- latent distance proportional to label
distance, $\rho |\Delta \tilde y|$. Both replace CEB's unconstrained reference with a
geometry that encodes the label space by construction, and both preserve the
posterior, the single KL, and the certificate.

\subsubsection{Positioning and extensions}
\label{sec:positioning}

The framework is built on CEB and inherits its residual certificate; the difference is
exclusively in the reference. CEB's backward encoder is an unconstrained trainable map
that can chase the posterior and merge classes --- the two failure channels of
Section~\ref{sec:theory}. The geometric prior is also trainable --- the frame rotates
and the axis turns --- but its degenerations are excluded by construction: it cannot
merge classes, collapse to a point, inflate its covariance, or rescale its geometry,
because orthonormality, scale, and covariance are fixed by the construction itself.
This is the sense in which the compression term measures distance to a \emph{declared}
target. The certificate of Eq.~\eqref{eq:opb-certificate} remains the usual CEB
residual bound for this particular prior family --- we do not claim a tighter mutual
information bound --- and the isometric axis likewise trades a general conditional
prior for a checkable label-metric structure. For multi-dimensional regression labels
$y \in \R^m$ ($m \le d$), the same construction extends verbatim: take
$W \in \R^{d \times m}$, its reduced QR $W = QR$, and the prior
$\mu_p(y) = \rho\, Q\, \tilde y$, which is isometric in the label metric; the scalar
case is $m = 1$ with $Q = u$. The label metric must be declared before
standardization, otherwise ``isometric'' has no unique meaning.

% =====================================================================
\section{Experiments}
\label{sec:experiments}
% =====================================================================

% 占位：后续填充实验部分。计划包含：
%   - 主扫参（baseline 对比，9 数据集 × 4 骨干）；
%   - 机制验证：先验几何分析（prior_geometry）、后验几何传递（posterior_geometry）、
%     最近锚点分类与能量分类器消融（geometry_usage + OPB-EC，§附录 prop:opb-energy）；
%   - 锚点尺度扫描（a ∈ {1,...,10}，半径跟随率 / Fisher / 跨 seed 一致性）；
%   - RandomLabel 记忆实验与对抗鲁棒性（adv_eval）。

% =====================================================================
\section{Related work}
\label{sec:related}
% =====================================================================

% 占位：后续填充相关工作（VIB / CEB / NIB / DVCCA / Caveats 已在前文引用，
% 此处补充神经坍缩、几何先验、分类头约束等邻近工作）。

% =====================================================================
\appendix
\section{Deferred proofs and further results}
\label{app:deferred}
% =====================================================================

The following results support the main text: the pointwise KL decomposition
(Lemma~\ref{lem:opb-kl-decomposition}), the anchor-energy classifier equivalence
(Proposition~\ref{prop:opb-energy}), a Fano-type label-information lower bound
(Corollary~\ref{cor:opb-fano}), the strict convexity of the $\beta$-sweep on the
aligned family (Proposition~\ref{prop:opb-sweep}), and the scope and limitations
(Remark~\ref{rem:opb-scope}).

\begin{lemma}[OPB KL decomposition]
\label{lem:opb-kl-decomposition}
For a diagonal posterior covariance and the isotropic prior of
Eq.~\eqref{eq:opb-prior},
\begin{equation}
\KL\bigl(q_\phi(z|x)\,\|\,p_\theta^{\mathrm{OPB}}(z|y)\bigr)
= \frac{1}{2\tau^2}\|\mu_\phi(x) - a\,q_{\theta,y}\|^2
+ \frac12 \sum_{r=1}^{d}\Bigl[
\frac{\sigma_{\phi,r}^2(x)}{\tau^2} - 1 - \log\frac{\sigma_{\phi,r}^2(x)}{\tau^2}
\Bigr],
\label{eq:opb-kl-decomposition}
\end{equation}
a class-anchor mismatch term plus a variance mismatch term that is non-negative and
vanishes exactly at $\sigma_{\phi,r}^2(x) = \tau^2$ for all $r$. The mismatch term is
quadratic in the displacement to a reference whose scale and covariance cannot move:
this is the mechanism behind the strict convexity of
Proposition~\ref{prop:opb-sweep} below.
\end{lemma}

\begin{proof}
Apply the diagonal-Gaussian KL formula with posterior mean $\mu_\phi(x)$, prior mean
$a\,q_{\theta,y}$, posterior variances $\sigma_{\phi,r}^2(x)$, and prior variances
$\tau^2$.
\end{proof}

\begin{proposition}[Anchor-energy classifier equivalence]
\label{prop:opb-energy}
Suppose the classifier uses the anchor energy
\begin{equation}
s_k(z) = -\frac{\|z - a\,q_{\theta,k}\|^2}{2\tau^2} + \log \pi_k .
\label{eq:opb-energy}
\end{equation}
Then $s_k(z)$ is the class-$k$ Gaussian log posterior up to a class-independent term;
equivalently, for uniform classes the equivalent linear classifier has weights
$w_k = (a/\tau^2)\,q_{\theta,k}$ --- equal norms and pairwise orthogonal rows. An
anchor-energy classifier therefore makes the prediction rule use exactly the same
geometry as the prior, and removes the free classifier scale by construction.
\end{proposition}

\begin{proof}
Expand the squared norm in Eq.~\eqref{eq:opb-energy}; the terms
$-\|z\|^2/(2\tau^2)$ and $-a^2/(2\tau^2)$ are common to all classes, and the
class-dependent remainder is the Gaussian discriminant score.
\end{proof}

\begin{corollary}[Label-information lower bound in the aligned Gaussian model]
\label{cor:opb-fano}
Assume uniform classes and exact alignment
$q(z|Y=k) = \mathcal{N}(a\,q_{\theta,k}, \tau^2 I_d)$. The pairwise error probability
of the nearest-anchor classifier is
$P_{\mathrm{pair}} = \Phi\bigl(-\frac{a}{\sqrt{2}\,\tau}\bigr)$, with
$P_e \le \min\{1, (K-1)P_{\mathrm{pair}}\}$, and whenever the displayed bound
$\bar P_e \le 1 - 1/K$, Fano's inequality gives
$I(Y;Z) \ge \log K - h_2(\bar P_e) - \bar P_e \log(K-1)$. This links the anchor
scale-to-noise ratio $a/\tau$ to the retained label information, conditionally on the
aligned Gaussian model.
\end{corollary}

\begin{proof}
Two anchors have distance $\sqrt{2}a$ by Proposition~\ref{prop:opb-frame}; the
equal-covariance binary decision boundary is halfway between them, giving
$\Phi(-\sqrt{2}a/(2\tau))$. The union bound and Fano's inequality on
$[0, 1-1/K]$ conclude.
\end{proof}

\begin{proposition}[Strictly convex $\beta$-sweep on the aligned family]
\label{prop:opb-sweep}
Consider the class-symmetric aligned family: every input of class $y$ is mapped to the
posterior mean $\mu_y = s\,q_{\theta,y}$ ($s \ge 0$), the classifier is held at
$w_j = c\,q_{\theta,j}$, $b_j = 0$ with fixed scale $c > 0$, and the posterior is
$q_y = \mathcal{N}(s\,q_{\theta,y}, \sigma^2 I)$. On the evaluation code ($z = \mu$),
\begin{equation}
\KL(s, \sigma^2) = \tfrac{d}{2}\bigl(\sigma^2/\tau^2 - 1 - \ln(\sigma^2/\tau^2)\bigr)
+ \tfrac{1}{2\tau^2}(s-a)^2,
\qquad
\CE(s) = \ln\bigl(1 + (K-1)e^{-cs}\bigr),
\label{eq:opb-kl-ce}
\end{equation}
and for every $\beta > 0$ the Lagrangian
$L_\beta(s) = \CE(s) + \tfrac{\beta}{2\tau^2}(s-a)^2$ is strictly convex in $s$ with a
unique minimizer $s^*(\beta) > a$ given by
$\tfrac{\beta}{\tau^2}(s-a) = c\,\varphi(cs)$ with
$\varphi(u) = (K-1)e^{-u}/(1+(K-1)e^{-u})$. The map $\beta \mapsto s^*$ is strictly
decreasing with $s^* \to a$ as $\beta \to \infty$ and $s^* \to \infty$ as
$\beta \to 0$, so $\beta \mapsto (\KL(s^*), \CE(s^*))$ is one-to-one onto the entire
trade-off curve: no two $\beta$ share an optimum, and every interior point is attained.
The fixed-scale condition is load-bearing: with $c$ free, the pair
$(s, c) = (a, c)$ with $c \to \infty$ drives both terms to zero and the infimum is $0$
for every $\beta$.
\end{proposition}

\begin{proof}
$\CE'(s) = -c\varphi(cs)$ and
$\CE''(s) = c^2(K-1)e^{-cs}/(1+(K-1)e^{-cs})^2 > 0$, and the compression term is
quadratic in $s-a$, so $L_\beta''(s) > 0$. The stationarity condition is exactly
$\tfrac{\beta}{\tau^2}(s-a) = c\varphi(cs)$; the left side increases from
$-\beta a/\tau^2$ through $0$ to $+\infty$ while the right side decreases within
$(0, c)$, giving a unique crossing with $s^* > a$. Differentiating with respect to
$\beta$ gives $s' = -(s-a)/(\beta - c^2\tau^2\varphi'(cs)) < 0$ since $\varphi' < 0$;
the limits follow by monotonicity. With $c$ free,
$\KL = \tfrac{1}{2\tau^2}(s-a)^2$ vanishes at $s = a$ while
$\CE = \ln(1+(K-1)e^{-ca}) \to 0$ as $c \to \infty$.
\end{proof}

\begin{remark}[The sweep is a KL--CE surrogate curve; scope and limitations]
\label{rem:opb-scope}
Proposition~\ref{prop:opb-sweep} proves strict convexity and one-to-one
parameterization in the KL--CE plane on the aligned family at fixed classifier scale,
not on the information plane of the Caveats paper; the results of the main text
establish a structured conditional prior, an exact KL geometry, and conditional
separation guarantees, and they do \emph{not} establish a mutual information upper
bound tighter than CEB's --- Eq.~\eqref{eq:opb-certificate} is the usual CEB gap
identity for a particular prior family. The QR construction fixes all pairwise anchor
distances, so the framework should not claim to recover arbitrary semantic class
similarities. The stateless construction is well-defined and locally differentiable
only while $M_\theta$ has full column rank (with a consistent sign convention); near
rank deficiency the frame can vary sharply. If the QR were computed from the current
input batch rather than from the all-class prior matrix, the prior would become
batch-dependent and Proposition~\ref{prop:opb-certificate} would no longer follow
directly. Finally, the free-scale caveat of Proposition~\ref{prop:opb-sweep} is
precisely what the anchor-energy classifier of Proposition~\ref{prop:opb-energy}
removes by construction --- the basis of the classifier ablation in the experiments.
\end{remark}

% =====================================================================
\bibliographystyle{plainnat}
\begin{thebibliography}{10}

\bibitem[Alemi et al.(2017)]{alemi2017deep}
A.~A.~Alemi, I.~Fischer, J.~V.~Dillon, and K.~Murphy.
\newblock Deep variational information bottleneck.
\newblock In \emph{ICLR}, 2017.

\bibitem[Fischer(2020)]{fischer2020conditional}
I.~Fischer.
\newblock The conditional entropy bottleneck.
\newblock \emph{Entropy}, 22(9):999, 2020.

\bibitem[Kolchinsky and Tracey(2019)]{kolchinsky2019caveats}
A.~Kolchinsky and B.~D.~Tracey.
\newblock Caveats for information bottleneck in deterministic scenarios.
\newblock In \emph{ICLR}, 2019.

\bibitem[Kolchinsky et al.(2017)]{kolchinsky2017nonlinear}
A.~Kolchinsky, B.~D.~Tracey, and D.~H.~Wolpert.
\newblock Nonlinear information bottleneck.
\newblock \emph{Entropy}, 21(12):1181, 2019.

\bibitem[Tishby et al.(2000)]{tishby2000information}
N.~Tishby, F.~C.~Pereira, and W.~Bialek.
\newblock The information bottleneck method.
\newblock In \emph{Proc.\ 37th Allerton Conference}, 2000.

\bibitem[Abdelaleem et al.(2025)]{abdelaleem2025deep}
E.~Abdelaleem, I.~Ndiour, and A.~Bhattacharyya.
\newblock Deep variational multivariate information bottleneck.
\newblock \emph{JMLR}, 2025.

\bibitem[Papyan et al.(2020)]{papyan2020prevalence}
V.~Papyan, X.~Y.~Han, and D.~L.~Donoho.
\newblock Prevalence of neural collapse during the terminal phase of deep learning
training.
\newblock \emph{PNAS}, 117(40):24652--24663, 2020.

\end{thebibliography}

\end{document}
